% Reproducibility sheet template. Filled by scripts/15_build_repro_sheet.py: every double-angle-bracket
% placeholder is computed from artifacts/ and reports/ at build time. Edit the prose here, never the numbers.
\documentclass[11pt]{article}
\usepackage[margin=1in]{geometry}
\usepackage{booktabs,longtable,array,graphicx,fancyvrb,enumitem,caption}
\usepackage{xurl}
\usepackage[hidelinks]{hyperref}
\setlength{\parindent}{0pt}
\setlength{\parskip}{6pt}
\setlength{\emergencystretch}{3em}
\captionsetup{font=small,labelfont=bf,skip=4pt}
\setlength{\LTpre}{6pt}
\setlength{\LTpost}{6pt}
\newcolumntype{P}[1]{>{\raggedright\arraybackslash}p{#1}}

\begin{document}

\begin{center}
{\LARGE\bfseries Reproducibility Sheet}\\[6pt]
{\large Multi-Turn Reinforcement Learning on Amazon SageMaker AI:\\
Training a Customer-Support Agent on GPT-OSS-20B}\\[12pt]
{\large Isham Rashik}\\[4pt]
Run of record: <<RUN_DATES>>
\end{center}

This sheet accompanies the workshop notebook \texttt{notebooks/sagemaker\_mtrl\_real\_demo.ipynb} and the
repository's \texttt{README.md}. With it, anyone with a clean copy of the repository can rebuild every table,
figure, and headline number from the recorded evidence without an AWS account, and can repeat the study in
their own AWS account. The interpretation lives in the notebook and the README; what follows is the recipe.
Every step is a numbered script in \texttt{scripts/} that imports the shared \texttt{aws} package, and the
agent that SageMaker trained against is in \texttt{agent/}. This sheet is itself generated by
\texttt{scripts/15\_build\_repro\_sheet.py}: every number below is read from \texttt{artifacts/} and
\texttt{reports/} when the sheet is built, and the build refuses to run unless every invariant passes, so the
sheet cannot drift from the evidence.

\section{Repository and Run Identifiers}
<<IDENTIFIERS_TABLE>>
<<IDENTIFIERS_NOTE>>

\section{Hardware and Operating System}
All model computation ran on AWS managed services. The local machine only submits and re-attaches to jobs,
drives the live inference, and builds the reports, so local hardware does not affect any reported number.
<<HARDWARE_TABLE>>

\section{Software Environment}
The local environment is a conda environment named \texttt{aws-mtrl} (Python 3.12). Three levels of pinning
are committed:
\begin{itemize}[leftmargin=*]
\item \texttt{environment-macos.yaml} and \texttt{requirements-macos.txt}: the direct dependencies at exact
versions (run of record: macOS). \texttt{environment-linux.yaml} and \texttt{requirements-linux.txt} carry the
same pins and have not been verified on Linux.
\item \texttt{requirements-macos.lock.txt}: every installed distribution (<<LOCK_COUNT>>) of the environment
that ran the scripts and the live inference.
\item \texttt{agent/requirements.lock.txt}: the exact <<AGENT_PACKAGES>> packages of the container image that
served every rollout, read from its CodeBuild log, with the base-image digest.
\end{itemize}
\begin{Verbatim}[frame=single,fontsize=\small]
conda env create -f environment-macos.yaml      # or environment-linux.yaml
conda activate aws-mtrl
pip install -r requirements-macos.lock.txt      # optional: the exact full environment
\end{Verbatim}
<<PACKAGE_TABLE>>

\section{Data}
Two committed prompt files define the study, and nothing is resampled. Training needs more prompts than the
global batch size (<<BATCH_SIZE>>), which is why the training set grew from 32 to 64 tickets before the
training job; the original 32 are kept in \texttt{artifacts/training\_prompts.32.csv}.
<<DATASET_TABLE>>
The held-out set never enters training: the two files share no prompt (invariant~<<INV_DISJOINT>>), and both
are byte-identical to the run of record (invariant~<<INV_FINGERPRINTS>>). The CSVs were converted to Parquet
and uploaded to S3 before the jobs ran. \texttt{scripts/11b\_snapshot\_provenance.py} later compared those S3
objects with the CSVs, prompt by prompt, and recorded when they were written
(invariant~<<INV_S3>>):
<<DATA_PROVENANCE_TABLE>>

\section{Random Seeds and Sources of Randomness}
<<SEEDS_TABLE>>
The study is therefore reproducible at two levels. Exactly: the offline steps rebuild every table, figure,
and this sheet from the evidence with identical values. Statistically: a new run uses the same code, data,
settings, and pins, but service-side sampling cannot be seeded, so its numbers differ within sampling noise.
With <<EVAL_ROLLOUTS>> rollouts per model, pass@1 has a standard error of about <<PASS1_SE>>. The same base
model scored pass@1 <<BASELINE_PASS1>> in the baseline and <<COMPARISON_BASE_PASS1>> in the comparison, on
the same tickets.

\section{Compute Accounting and Budgets}
Costs are recorded per job and never silently merged. Token-priced lines use the token counts AWS billed
(\texttt{BillableTokenUsage}) at the AWS Price List rates in effect, in USD per million tokens: <<RATES>>.
The other lines are computed or estimated, and their labels say how. AWS Cost Explorer is the final bill.
Each line's token usage is in \texttt{artifacts/cost\_accounting.json}; the table is
\texttt{reports/repro/compute\_accounting.csv}.
<<COST_TABLE>>
<<COST_TOTALS>>
<<DURATION_TABLE>>
A budget guard (\texttt{aws/costs/cost\_guard.py}) runs before training, the comparison evaluation, the
endpoint, the Bedrock import, and live inference. It adds the spend so far to a projection of the planned
steps and refuses any plan above the \$<<BUDGET_CAP>> cap. The projection scales this account's billed
baseline tokens per rollout to the planned rollouts with a 1.5$\times$ margin, and adds the endpoint's
maximum lifetime at the official hourly price of \$<<HOSTING_RATE>>.

\section{Experimental Protocol and Disclosures}
These are the methodology choices and deviations a reviewer most needs to see stated explicitly. The
machine-readable record is \texttt{reports/repro/study\_metadata.json}.
<<PROTOCOL_TABLE>>

\section{Method Configuration Disclosures}
Every setting below is fixed in \texttt{aws/config.py}, pinned as the run-of-record protocol in
\texttt{aws/reporting/invariants.py}, or recorded by the AWS job itself.
<<CONFIG_TABLE>>

\section{Metric-Calculation Conventions}
The evaluator reports its metrics by name; a reproduction needs the exact definition. Each definition marked
``recomputed'' was checked by recomputing the reported value from the recorded conversations when this sheet
was built, and the solved count, pass@1, and mean reward are also checked by invariant~<<INV_TRACES>>.
<<METRICS_TABLE>>

\section{Reproducing Each Figure and Table}
From a clean checkout, create the environment (Section~3) and run the numbered scripts in order.
\texttt{reports/repro/run\_commands.csv} is the authoritative, literal command sequence. A billable step runs
only with its exact confirmation phrase; without the phrase a script shows the recorded result, and a recorded
job is re-attached, never submitted twice.

\textbf{A. Rebuild every reported number offline} (free; no AWS account needed):
<<COMMANDS_OFFLINE>>
\textbf{B. Re-attach to the recorded jobs} (the run's account; read-only, submits nothing):
<<COMMANDS_REATTACH>>
\textbf{C. Repeat the study in your own account} (about \$6--8). First move this run's evidence aside
(\texttt{mv artifacts artifacts-run-of-record}, \texttt{mv reports reports-run-of-record},
\texttt{mkdir artifacts}), then:
<<COMMANDS_FULL>>
The notebooks present the same results with narrative; open them with
\texttt{python -m jupyter lab notebooks/} from the project root.
<<NOTEBOOKS_TABLE>>

\section{Headline Results (for a quick check)}
An offline rebuild lands exactly on the numbers below. A new run should land within sampling noise, with the
fine-tuned model ahead of the base model on pass@1.
<<HEADLINE_TABLE>>

\section{Extended Materials}
Tables are from \texttt{reports/tables/} and figures from \texttt{reports/figures/}. Both are rebuilt by
\texttt{scripts/12\_make\_report\_tables\_and\_figures.py}.

\subsection{Training}
<<FIG_TRAINING>>
<<TABLE_TRAINING>>

\subsection{Held-out evaluation}
<<FIG_EVALUATION>>
<<TABLE_EVALUATION>>
<<TABLE_BASELINES>>

\subsection{Per held-out ticket}
<<FIG_TICKETS>>
<<TABLE_TICKETS>>

\subsection{Live inference}
<<FIG_LIVE>>
<<TABLE_LIVE>>

\subsection{Run of record}
<<TABLE_RUN_OF_RECORD>>

\subsection{Provenance: what actually ran}
<<PROVENANCE_TEXT>>
<<TABLE_PROVENANCE>>
<<TABLE_AGENT_FILES>>
<<TABLE_IMPORT_FILES>>

\section{Invariant Checks}
\texttt{scripts/14\_verify\_invariants.py} re-checks, offline and from the files on disk, every property the
results depend on, and prints a PASS/FAIL line for each. <<INVARIANT_SUMMARY>>
<<INVARIANT_LIST>>

\section{Held-out Discipline and Provenance Controls}
The <<EVAL_PROMPTS>> held-out tickets are fixed before any training and never enter it. The two prompt files
share no prompt (invariant~<<INV_DISJOINT>>), and both match the run of record by SHA-256
(invariant~<<INV_FINGERPRINTS>>). Training reads only the training tickets. The baseline, the comparison, and
the live inference read only held-out tickets. The live inference uses the first <<LIVE_TICKETS>> held-out
tickets in file order with the documented seed (invariant~<<INV_LIVE>>). Both models in the comparison are
evaluated by one pipeline, on the same tickets, through the same agent and serving stack, so the comparison
isolates the effect of training.

Provenance closes the remaining gap between the repository and what ran:
\begin{itemize}[leftmargin=*]
\item \textbf{Agent.} The AgentCore runtime that served every rollout went live before all three jobs
started, and was built from the saved source bundle. That bundle's environment, rollout driver, Dockerfile,
and system prompt match the repository (invariant~<<INV_AGENT>>).
\item \textbf{Datasets.} The S3 datasets the jobs read equal the CSVs, and were written before the jobs
(invariant~<<INV_S3>>).
\item \textbf{Imported model.} It was built from the trained package's merged weights, changing only the
two files the Bedrock documentation requires (invariant~<<INV_DEPLOY>>).
\end{itemize}

\appendix
\section{Software Engineering Practices}
The study code is treated as software, not throwaway scripts. Each practice below names the concrete
mechanism in the codebase.
\begin{itemize}[leftmargin=*]
\item \textbf{Separation of concerns.} The \texttt{aws} package is organised by stage: \texttt{setup/},
\texttt{rl/}, \texttt{deploy/}, \texttt{costs/}, \texttt{records/}, and \texttt{reporting/}, with
\texttt{workflow.py} as the single facade. The numbered scripts are thin entry points, and \texttt{agent/} is
kept identical to what ran.
\item \textbf{DRY (single source of truth).} Settings and project paths live once in \texttt{aws/config.py}.
One table builder serves both script 12 and the invariants. The notebook and the scripts call the same
modules, and the notebook is generated by \texttt{scripts/dev/build\_notebook.py} and checked against it.
\item \textbf{Immutable values and write-once evidence.} Settings and the protocol are frozen dataclasses.
The evidence of record is never overwritten: a later live run is saved beside it, and a failed job is
archived before it is replaced.
\item \textbf{Fail-fast errors.} A missing evidence file raises \texttt{FileNotFoundError} naming the script
that produces it. A wrong confirmation phrase raises \texttt{PermissionError}. The budget guard refuses before
any billable work.
\item \textbf{Cost safety.} Every billable step needs an exact confirmation phrase, passes the budget guard,
and re-attaches to a recorded job instead of resubmitting it. A watchdog deletes the endpoint at its
deadline.
\item \textbf{Reproducibility hygiene.} The project uses only project-relative paths, and the environments
are pinned and locked. \texttt{run\_commands.csv} is ordered, \texttt{study\_metadata.json} is
machine-readable, and evidence is written atomically with signed links redacted. Logs are scrubbed of links
and of the home-folder path.
\item \textbf{Verify before asserting.} The invariants recompute the metrics from the recorded conversations,
the costs from token usage and rates, and the tables from the evidence. This sheet is built only when all of
them pass.
\item \textbf{Tests and style.} <<TEST_COUNT>> unit tests (\texttt{tests/} mirrors \texttt{aws/};
\texttt{agent/} has its own) run with \texttt{python -m pytest} and need no AWS access. The code has
Google-style docstrings and type hints throughout, and passes \texttt{ruff}.
\end{itemize}

\end{document}
